\section{Learning with finite trusted control specifications}
\label{sec:finitecontrol78}

The learning theorems count unknown-device calls.  We now remove
exact physical preparation from their hypotheses, while keeping
their loss and call orders.  A trusted preparation instruction has a
finite classical description and a certified trace-norm error; a
trusted control instruction has a finite description and a certified
diamond-norm error.  These are accuracy assumptions on the available
controls.  They assert neither a universal hardware implementation
nor a bound on preparation time, gate count, or classical workspace.
All trace norms in this section are unhalved.

\begin{lemma}[An adaptive control-error budget]
\label{lem:controlhybrid78}
Fix a finite classical decision tree of experiments, with bounded
numbers of calls and trusted instructions on every record.  Its
unknown channel is unchanged.  At a scheduled trusted event $j$,
conditional on the preceding classical history $h$, replace a
channel $\mathcal C_{j,h}$ by a channel $\widetilde{\mathcal C}_{j,h}$
such that
\begin{equation}\label{eq:controlbudget78}
 \|\widetilde{\mathcal C}_{j,h}-\mathcal C_{j,h}\|_\diamond
 \le e_j(h),\qquad
 \sum_{j\text{ on }h_*}e_j(h_{<j})\le e
 \quad\hbox{for every complete record }h_*.
\end{equation}
For a fresh state preparation, its trace-norm error suffices in
place of the diamond norm: conditional on $h$, the new state is
tensor-separated from every carried system in both experiments.
The final classical--quantum states
satisfy
\begin{equation}\label{eq:controlhybrid78}
 \|\widetilde\rho_{\rm out}-\rho_{\rm out}\|_1\le e.
\end{equation}
In particular, the laws of any unchanged classical output have
total variation distance at most $e/2$.
\end{lemma}
\begin{proof}
Include each observed outcome in the classical history and pad
stopped branches by identities.  Quantum registers of different
sizes can be placed in a finite direct sum.  Each event is then a
channel on the complete classical--quantum state.  A quantum
instrument means its complete channel
$\rho\mapsto\sum_x|x\rangle\langle x|\otimes\mathcal C_x(\rho)$:
if errors are instead supplied separately for its completely
positive branches, their sum is the event's error charge.

At each event, insert the ideal preceding state between the actual
and ideal updates.  Contractivity of the actual channel and the
diamond-norm bound give
\[
 D_j\le D_{j-1}
       +\sum_h p_{j-1}(h)e_j(h),
\]
where $D_j$ is the unhalved trace-norm difference and $p_{j-1}$ is
the ideal history law.  Unknown-channel calls and unchanged
classical calculations have charge zero.  Summing the inequality
gives the ideal expectation of the pathwise sum in
\eqref{eq:controlbudget78}, which is at most $e$.
A replacement preparation channel has diamond distance equal to
the trace distance of the prepared states, including when another
register is retained.  Finally, classical readout and deterministic
postprocessing are contractions, and classical total variation is
one half of the trace norm.  No continuity of the decision rules
or of the final loss is used.
\end{proof}

We check that the learners to which this lemma is applied have
finite, computable control specifications.  Take integer dimensions
and horizons and rational positive accuracy and confidence targets.
All absolute proof constants may be fixed once as sufficiently
conservative rational numbers.  For sample budgets, replace each
$\log x$, $x>1$ rational, by
\[
 \ell(x)=\min\{k\in\mathbb N:2^k\ge x\}.
\]
Here $\log x\le\ell(x)\le1+\log x/\log2$.
Thus the concentration bounds and call orders are preserved, while
ceilings of the resulting rational sample budgets are decidable by
integer arithmetic.  This convention avoids requiring an exact
comparison of a transcendental expression with an integer.

Calibration frequencies and all complex empirical means are
rational.  The ordered spectral projections, fixed-order
Gram--Schmidt procedure, clipping, normalized frames, and tie rules
of Sections~\ref{sec:commonlearn76}--\ref{sec:matrixlearning77} use
algebraic arithmetic.  The notation for phase recovery does not
require a transcendental oracle.  For a nonzero empirical $z$ passing
\eqref{eq:learnphaseguard77}, put $q=z/|z|$ and select the root $v$
of $v^m=q$ with largest real part.  It is unique: the passed guard
places $\operatorname{Arg}q$ strictly between $-\pi$ and $\pi$,
so its principal $m$th root is strictly closer in angle to $1$
than every other root.  All roots are algebraic, and comparison of
their real parts is decidable.  Consequently
\[
 \tan\bigl(\operatorname{Arg}(z)/m\bigr)
       =\frac{\Im v}{\Re v}
\]
is algebraic; its denominator is positive under the guard.
Zero, amplitude, frame-selection, and phase-guard equalities are
therefore decidable exactly.  The finite rational legalization grid
in Theorem~\ref{thm:matrixlearning77} is in the public input basis;
its positivity tests are exact, and its finitely many algebraic
objectives can be evaluated with the prescribed slack.  The same
holds for the exact rational codec.  Every record consequently has
a finite, terminating classical computation and a legal output.

There are finitely many records at fixed public parameters: the
device outputs and preparation signs have finite alphabets and
their numbers are bounded on every path.  This does not discretize
the unknown effect or the retained quantum state.  Within a block
the state still lies in its full continuous quantum state space.
Only the experiment's classical instructions and observations have
finite descriptions.

Finite descriptions can also specify controls with any required
positive accuracy.  If a computable unit vector $\psi\in\mathbb C^D$
is prescribed and $0<\epsilon\le1$, compute a nonzero Gaussian
rational vector $z$ with
$\|z-\psi\|_2\le\epsilon/8$ and set
$\phi=z/\sqrt{\sum_k|z_k|^2}$.  Then
\begin{equation}\label{eq:controlstate78}
 \bigl\||\phi\rangle\langle\phi|
              -|\psi\rangle\langle\psi|\bigr\|_1
 \le2\|\phi-\psi\|_2\le4\|z-\psi\|_2\le\epsilon/2.
\end{equation}
Coordinate accuracy $O(\epsilon/\sqrt D)$ suffices, so $z$ and
its normalization rule have a description of
$O(D\log(D/\epsilon))$ bits.  The normalization is specified
algebraically; rational normalized amplitudes are not assumed.
A trusted preparation within $\epsilon/2$ of this finitely
specified target has total error at most $\epsilon$ from the ideal
state.

For completeness, finite specifications can preserve the exact
channel constraints as well.  Approximate a known algebraic
Stinespring isometry $V$ by a Gaussian rational matrix $A$ with
$\|A-V\|_{\rm op}\le r<1/2$, and specify
$W=A(A^*A)^{-1/2}$.  Then $W^*W=I$ and
$\|W-V\|_{\rm op}\le2r$: every singular value of $A$ differs
from one by at most $r$.  The channels obtained from these
isometries differ in diamond norm by at most $4r$, by expansion of
$W\rho W^*-V\rho V^*$ and partial-trace contractivity.  Keeping a
finite outcome flag gives the same construction for instruments.
Their actual trusted realizations are charged separately in
\eqref{eq:controlbudget78}.  Independent rounding of Kraus entries
without restoring the channel constraints would not suffice.

\begin{theorem}[Finite-control form of the learning guarantees]
\label{thm:finitecontrollearning78}
The common matrix learner of Theorem~\ref{thm:matrixlearning77}
and the learner with training-block cap $1\le b\le N$ of
Theorem~\ref{thm:blocklearning78} admit the preceding finite
trusted-control model in their stated ranges.
They have the original success target
\[
 \sup_{E\in\mathfrak E_d}
 \mathbb P_E\{d_N(\mathcal M_E,\mathcal M_{\widehat E})>\delta\}
 \le\eta
\]
with every-record call orders, respectively,
\begin{equation}\label{eq:finitecontrolcalls78}
 Cd^4N\delta^{-2}\log(d/\eta),\qquad
 Cd^4\frac{N^2}{b}\delta^{-2}\log(d/\eta).
\end{equation}
Their original block caps and exact legality of the returned
classical effect are preserved.  The same transfer applies when the
classical output is the public codeword.
In dimension one the order remains
$CN\delta^{-2}\log(1/\eta)$ independently of $b$.
\end{theorem}
\begin{proof}
Use the corresponding ideal learner at the same accuracy $\delta$
and confidence $\eta/2$, with the finite sample-budget convention
above.  Fix a computable public upper bound $M_*$ on its calls
on every record.  The explicit constructions use fresh product
states or fresh GHZ states embedded into a known two-dimensional
subspace.  They can be implemented by one preparation event per
block, hence at most $M_*$ such events.  Given any preceding
record, compute that record's prescribed state to trace accuracy
$\eta/M_*$, using \eqref{eq:controlstate78} and the trusted
realization allowance.  Lemma~\ref{lem:controlhybrid78} gives final
output total variation at most $\eta/2$.  For each fixed unknown
$E$, apply it to the set of returned effects whose original loss
exceeds $\delta$.  The ideal probability is at most $\eta/2$,
so the actual probability is at most $\eta$.  Discontinuous
decisions, unsuccessful records, and exact equality branches are
included in this comparison.

Classical postprocessing is identical in the two experiments; it
does not round the returned effect or test a perturbed loss.
Thus legality is exact on every actual record and no additional
$\delta$-dependent loss allowance is needed.  Halving $\eta$ and
upper-rounding sample budgets change only the absolute constants
in \eqref{eq:finitecontrolcalls78}.
\end{proof}

Each approximate preparation above supplies a state on the same
$m\le b$ probe registers as its ideal block; preparation workspace
is discarded or reset before the experiment and no quantum system
is passed to the next block.  Approximation therefore does not
enlarge entangled input width or introduce coherent coupling from
old memory into fresh probes.  In a realization with $S$ additional
scheduled control instructions, use a public bound on the total
number of trusted events and charge every instruction in
\eqref{eq:controlbudget78}.  Per-event norm precision is then of
order $\eta/(M_*+S)$.  For an $m$-probe state the direct coordinate
description has $D=d^m$ and size
$O(D\log(D(M_*+S)/\eta))$ bits.  The precision per coordinate
is logarithmic in $(M_*+S)/\eta$, with its dimension dependence
explicitly charged; this is not a polynomial preparation bound.
Real-valued target tolerances may be replaced by supplied rational
targets between one half and the requested values, preserving all
orders and the requested guarantees.  No representation of an
arbitrary real input, and no oracle for the unknown effect, is
assumed.
